%  --- Définitions de chapitres ---
\newcolumntype{C}[1]{>{\centering\arraybackslash}m{#1}}
\newcolumntype{L}[1]{>{\raggedright\arraybackslash}m{#1}}
\newcommand{\protocolcols}{C{0.05\linewidth} L{0.26\linewidth} L{0.25\linewidth} C{0.13\linewidth} L{0.21\linewidth}}
\newcommand{\phdr}{%
    \textbf{N\textsuperscript{o}} & \textbf{Point(s) à mesurer} & \textbf{Valeur attendue} &
    \textbf{Valeur mesurée} & \textbf{OK / NOK + remarques} \\}
% ---------------------------------

\chapter{Protocole de test et de validation matérielle}
\label{chap:protocole-test}

\setlength{\tabcolsep}{4pt}

Cette annexe regroupe l'ensemble des mesures de mise en service (\textit{bring-up})
de la carte YDSTB2026. Les essais sont ordonnés de façon à valider progressivement
le matériel : contrôles hors tension, puis alimentations, puis communication et
programmation, puis périphériques numériques, et enfin la chaîne audio et les essais
fonctionnels. Cet ordre n'est pas arbitraire : on ne sollicite un étage que lorsque
son alimentation a été validée, afin de ne pas endommager un composant sur un rail
mal réglé ou court-circuité.

\paragraph{Précautions.}
La première mise sous tension se fait de préférence sur une alimentation de laboratoire
à \textbf{courant limité} (limite basse, p.~ex. 100--200~mA), afin qu'un court-circuit
non détecté à l'ohmmètre se traduise par une limitation de courant et non par une
destruction. L'\textbf{offset DC en sortie de jack} (mesure n\textsuperscript{o}~71)
est vérifié \emph{avant} de connecter un casque de valeur. Les niveaux audio sont
montés progressivement, en partant d'un volume numérique bas.

\paragraph{Matériel recommandé.}
Multimètre, alimentation de laboratoire à courant limité, oscilloscope, générateur de
signal ou fichier MP3 de test (1~kHz), charge audio (casque de test ou résistances),
et, en option, un analyseur audio pour le THD+N et le SNR.

\paragraph{Lecture des tableaux.}
Pour les contrôles visuels, la « valeur attendue » est l'état conforme recherché.
La colonne « Valeur mesurée » est laissée vide pour le relevé ; la dernière colonne
sert au verdict (OK/NOK) et aux remarques éventuelles.

% =====================================================================
\section{Contrôles avant mise sous tension}
% =====================================================================
\begin{longtable}{\protocolcols}
    \caption{Inspection visuelle et contrôles à l'ohmmètre (carte hors tension).}
    \label{tab:proto-preflight}                                                                                                                                                        \\
    \toprule \phdr \midrule \endfirsthead
    \multicolumn{5}{l}{\footnotesize\itshape Tableau \thetable{} — suite}                                                                                                              \\
    \toprule \phdr \midrule \endhead
    \midrule \multicolumn{5}{r}{\footnotesize\itshape suite page suivante}                                                                                                             \\ \endfoot
    \bottomrule \endlastfoot
    \rh 1  & Inspection visuelle générale (soudures, ponts, composants manquants ou décalés)                    & Aucun pont ni manque, soudures conformes                        &  & \\
    \rh 2  & Polarité des composants polarisés (TVS/diodes, condensateurs, connecteur batterie, U1, DAC, ampli) & Repères pin 1 / cathodes corrects                               &  & \\
    \rh 3  & Connecteur USB-C (alignement, absence de pont entre pads)                                          & Conforme                                                        &  & \\
    \rh 4  & Continuité du plan de masse entre points de test GND                                               & $\approx$ 0~$\Omega$                                            &  & \\
    \rh 5  & Court-circuit VBUS $\rightarrow$ GND                                                               & Pas de court (R non nulle, montée par charge des condensateurs) &  & \\
    \rh 6  & Court-circuit VSYS (sortie chargeur) $\rightarrow$ GND                                             & Pas de court                                                    &  & \\
    \rh 7  & Court-circuit VBAT $\rightarrow$ GND                                                               & Pas de court (si batterie connectée, on lit VBAT)               &  & \\
    \rh 8  & Court-circuit 3V3 (sortie TPS62A01) $\rightarrow$ GND                                              & Pas de court                                                    &  & \\
    \rh 9  & Court-circuit AVDD / CPVDD / DVDD du PCM5242 $\rightarrow$ GND                                     & Pas de court                                                    &  & \\
    \rh 10 & Court-circuit PVDD du MAX97220 $\rightarrow$ GND                                                   & Pas de court                                                    &  & \\
    \rh 11 & Isolement entre rails (VBUS / VSYS / 3V3 / VBAT entre eux)                                         & Pas de court inter-rails                                        &  & \\
    \rh 12 & Orientation du connecteur batterie (pas d'inversion +/$-$)                                         & Correcte                                                        &  & \\
\end{longtable}

% =====================================================================
\section{Mise sous tension et alimentations}
% =====================================================================
\begin{longtable}{\protocolcols}
    \caption{Alimentations et gestion d'énergie (sous tension, alimentation à courant limité).}
    \label{tab:proto-power}                                                                                                                                                                                 \\
    \toprule \phdr \midrule \endfirsthead
    \multicolumn{5}{l}{\footnotesize\itshape Tableau \thetable{} — suite}                                                                                                                                   \\
    \toprule \phdr \midrule \endhead
    \midrule \multicolumn{5}{r}{\footnotesize\itshape suite page suivante}                                                                                                                                  \\ \endfoot
    \bottomrule \endlastfoot
    \rh 13 & Courant d'appel à la 1\textsuperscript{re} mise sous tension (alim. limitée) & Pas de surintensité ; consommation au repos faible (à documenter)                                          &  & \\
    \rh 14 & VBUS au connecteur USB-C                                                     & 5,0~V (4,75--5,25~V)                                                                                       &  & \\
    \rh 15 & Entrée chargeur MAX77757 (CHGIN)                                             & $\approx$ VBUS (5~V), alimentation détectée valide                                                         &  & \\
    \rh 16 & Sortie système du chargeur (SYS / VSYS)                                      & Tension présente, cohérente avec le mode (à confirmer selon config.)                                       &  & \\
    \rh 17 & Tension batterie VBAT (cellule Li-ion)                                       & 3,0--4,2~V selon l'état de charge                                                                          &  & \\
    \rh 18 & Courant de charge (batterie partiellement déchargée)                         & Conforme à I\textsubscript{CHG} programmé (à confirmer registre MAX77757)                                  &  & \\
    \rh 19 & Tension de fin de charge (CV)                                                & 4,20~V ($\pm$50~mV) selon config.                                                                          &  & \\
    \rh 20 & Sortie buck TPS62A01 (rail 3V3)                                              & 3,30~V ($\pm$3~\%, soit 3,20--3,40~V)                                                                      &  & \\
    \rh 21 & Ondulation sur 3V3 (oscilloscope AC, en charge)                              & Faible (qq dizaines de mV cc max, à documenter)                                                            &  & \\
    \rh 22 & PCM5242 — AVDD                                                               & 3,3~V ($\pm$5~\%)                                                                                          &  & \\
    \rh 23 & PCM5242 — CPVDD (alim. charge pump)                                          & 3,3~V                                                                                                      &  & \\
    \rh 24 & PCM5242 — VNEG (rail négatif charge pump)                                    & $\approx$ $-$3,3~V (charge pump active)                                                                    &  & \\
    \rh 25 & PCM5242 — DVDD / sortie LDO interne                                          & Selon config. ($\approx$ 1,8~V si LDO interne, à confirmer datasheet)                                      &  & \\
    \rh 26 & MAX97220 — PVDD                                                              & 3,3~V                                                                                                      &  & \\
    \rh 27 & Superviseur XC6120N — sortie au repos (VBAT nominal)                         & Inactive (pas de mise en défaut)                                                                           &  & \\
    \rh 28 & Seuil de détection sous-tension XC6120N (abaisser VBAT via alim. labo)       & Bascule au seuil V\textsubscript{DF} de la référence ($\approx$ 3,0~V typ., selon suffixe) ; coupe/signale &  & \\
\end{longtable}

% =====================================================================
\section{Démarrage, EN et strapping}
% =====================================================================
\begin{longtable}{\protocolcols}
    \caption{Reset, strapping et démarrage du module ESP32-WROVER-E.}
    \label{tab:proto-boot}                                                                                              \\
    \toprule \phdr \midrule \endfirsthead
    \multicolumn{5}{l}{\footnotesize\itshape Tableau \thetable{} — suite}                                               \\
    \toprule \phdr \midrule \endhead
    \midrule \multicolumn{5}{r}{\footnotesize\itshape suite page suivante}                                              \\ \endfoot
    \bottomrule \endlastfoot
    \rh 29 & Broche EN après mise sous tension        & 3,3~V (relâchée après le RC)                               &  & \\
    \rh 30 & Broche GPIO0 au repos (mode boot normal) & Niveau haut (3,3~V)                                        &  & \\
    \rh 31 & Détection PSRAM au boot (log ESP-IDF)    & SPIRAM détectée, taille attendue (8~Mo)                    &  & \\
    \rh 32 & Log de boot lisible sur UART             & Boot ESP-IDF normal, pas de boucle de reset ni de brownout &  & \\
\end{longtable}

% =====================================================================
\section{USB, pont UART et programmation}
% =====================================================================
\begin{longtable}{\protocolcols}
    \caption{Pont CP2102N, multiplexeur TC7USB40MU et flash.}
    \label{tab:proto-usb}                                                                                                                            \\
    \toprule \phdr \midrule \endfirsthead
    \multicolumn{5}{l}{\footnotesize\itshape Tableau \thetable{} — suite}                                                                            \\
    \toprule \phdr \midrule \endhead
    \midrule \multicolumn{5}{r}{\footnotesize\itshape suite page suivante}                                                                           \\ \endfoot
    \bottomrule \endlastfoot
    \rh 33 & Énumération du CP2102N sur l'hôte                              & Port COM/tty détecté (Silicon Labs CP210x)                        &  & \\
    \rh 34 & Routage du mux TC7USB40MU (CP2102N $\leftrightarrow$ chargeur) & Commute selon la sélection ; D+/D$-$ vers CP2102N en mode debug   &  & \\
    \rh 35 & Auto-flash (esptool pilotant DTR/RTS)                          & Entrée en mode bootloader sans appui manuel ; EN et GPIO0 pilotés &  & \\
    \rh 36 & Connexion esptool (chip\_id / MAC)                             & Identifie ESP32, lit la MAC sans erreur                           &  & \\
    \rh 37 & Flash d'un firmware de test puis boot                          & Flash OK, redémarrage et exécution                                &  & \\
\end{longtable}

% =====================================================================
\section{Bus I\textsuperscript{2}C — détection des périphériques}
% =====================================================================
\begin{longtable}{\protocolcols}
    \caption{Scan I\textsuperscript{2}C (adresses à confirmer selon le câblage des broches d'adresse).}
    \label{tab:proto-i2c}                                                                                                           \\
    \toprule \phdr \midrule \endfirsthead
    \multicolumn{5}{l}{\footnotesize\itshape Tableau \thetable{} — suite}                                                           \\
    \toprule \phdr \midrule \endhead
    \midrule \multicolumn{5}{r}{\footnotesize\itshape suite page suivante}                                                          \\ \endfoot
    \bottomrule \endlastfoot
    \rh 38 & Détection FXL6408D                   & ACK à l'adresse attendue ($\approx$ 0x43, selon broche ADDR — à confirmer) &  & \\
    \rh 39 & Détection PCM5242                    & ACK (0x4C par défaut ; 0x4C--0x4F selon broches ADR)                       &  & \\
    \rh 40 & Détection MAX17260                   & ACK à 0x36                                                                 &  & \\
    \rh 41 & Détection MAX77757                   & ACK à l'adresse attendue (à confirmer datasheet)                           &  & \\
    \rh 42 & Absence d'adresse parasite / conflit & Seules les 4 adresses attendues répondent                                  &  & \\
\end{longtable}

% =====================================================================
\section{Expanseur GPIO FXL6408D}
% =====================================================================
\begin{longtable}{\protocolcols}
    \caption{Contrôle des sorties (mute DAC, shutdown ampli) et lecture des boutons.}
    \label{tab:proto-fxl}                                                                                                                    \\
    \toprule \phdr \midrule \endfirsthead
    \multicolumn{5}{l}{\footnotesize\itshape Tableau \thetable{} — suite}                                                                    \\
    \toprule \phdr \midrule \endhead
    \midrule \multicolumn{5}{r}{\footnotesize\itshape suite page suivante}                                                                   \\ \endfoot
    \bottomrule \endlastfoot
    \rh 43 & Sortie « mute DAC » $\rightarrow$ broche XSMT du PCM5242        & Niveau conforme à la commande (mute = actif)             &  & \\
    \rh 44 & Sortie « shutdown ampli » $\rightarrow$ broche SHDN du MAX97220 & Niveau conforme (shutdown = ampli coupé)                 &  & \\
    \rh 45 & Lecture des boutons (chacun : repos puis appui)                 & Le niveau logique change à l'appui, sur tous les boutons &  & \\
    \rh 46 & Interruption INT du FXL6408D vers l'ESP32                       & INT s'active au changement d'entrée ; ISR déclenchée     &  & \\
\end{longtable}

% =====================================================================
\section{Jauge de carburant MAX17260}
% =====================================================================
\begin{longtable}{\protocolcols}
    \caption{Lecture des registres de la jauge (ModelGauge m5).}
    \label{tab:proto-gauge}                                                                                    \\
    \toprule \phdr \midrule \endfirsthead
    \multicolumn{5}{l}{\footnotesize\itshape Tableau \thetable{} — suite}                                      \\
    \toprule \phdr \midrule \endhead
    \midrule \multicolumn{5}{r}{\footnotesize\itshape suite page suivante}                                     \\ \endfoot
    \bottomrule \endlastfoot
    \rh 47 & VCell (09h) vs VBAT mesuré au multimètre & Écart $<$ $\pm$20~mV                              &  & \\
    \rh 48 & RepSOC (06h)                             & 0--100~\%, plausible vis-à-vis de la tension      &  & \\
    \rh 49 & Current (0Ah)                            & Signe et amplitude cohérents (charge vs décharge) &  & \\
    \rh 50 & Temp (08h)                               & Température plausible (proche de l'ambiante)      &  & \\
    \rh 51 & Suivi du SOC pendant charge/décharge     & Le SOC évolue dans le bon sens                    &  & \\
\end{longtable}

% =====================================================================
\section{Chargeur MAX77757 (registres)}
% =====================================================================
\begin{longtable}{\protocolcols}
    \caption{État de charge et limites programmées.}
    \label{tab:proto-charger}                                                                                                                              \\
    \toprule \phdr \midrule \endfirsthead
    \multicolumn{5}{l}{\footnotesize\itshape Tableau \thetable{} — suite}                                                                                  \\
    \toprule \phdr \midrule \endhead
    \midrule \multicolumn{5}{r}{\footnotesize\itshape suite page suivante}                                                                                 \\ \endfoot
    \bottomrule \endlastfoot
    \rh 52 & Détection présence CHGIN (registre d'état)                 & « Input valid » quand l'USB est branché                                     &  & \\
    \rh 53 & État de charge (off / CC / CV / done)                      & Cohérent avec la réalité (la complétion vient du chargeur, pas de la jauge) &  & \\
    \rh 54 & Limite de courant d'entrée et courant de charge programmés & Conformes aux valeurs écrites                                               &  & \\
    \rh 55 & Absence de défaut en fonctionnement normal                 & Aucun flag (thermique, timer, etc.)                                         &  & \\
\end{longtable}

% =====================================================================
\section{Potentiomètre de volume (ADC)}
% =====================================================================
\begin{longtable}{\protocolcols}
    \caption{Lecture ADC du potentiomètre de volume.}
    \label{tab:proto-pot}                                                                                                               \\
    \toprule \phdr \midrule \endfirsthead
    \multicolumn{5}{l}{\footnotesize\itshape Tableau \thetable{} — suite}                                                               \\
    \toprule \phdr \midrule \endhead
    \midrule \multicolumn{5}{r}{\footnotesize\itshape suite page suivante}                                                              \\ \endfoot
    \bottomrule \endlastfoot
    \rh 56 & Lecture ADC, potentiomètre au minimum & Valeur brute proche de 0                                                      &  & \\
    \rh 57 & Lecture ADC, potentiomètre au maximum & Proche de la pleine échelle ($\approx$ 4095, 12~bits)                         &  & \\
    \rh 58 & Balayage sur toute la course          & Monotone, sans saut ni zone morte (tenir compte de la non-linéarité de l'ADC) &  & \\
\end{longtable}

% =====================================================================
\section{Affichage OLED SSD1333 (SPI2)}
% =====================================================================
\begin{longtable}{\protocolcols}
    \caption{Initialisation et rendu de l'écran.}
    \label{tab:proto-oled}                                                                                                                             \\
    \toprule \phdr \midrule \endfirsthead
    \multicolumn{5}{l}{\footnotesize\itshape Tableau \thetable{} — suite}                                                                              \\
    \toprule \phdr \midrule \endhead
    \midrule \multicolumn{5}{r}{\footnotesize\itshape suite page suivante}                                                                             \\ \endfoot
    \bottomrule \endlastfoot
    \rh 59 & Initialisation du SSD1333 (reset, séquence d'init) & Init sans erreur, écran allumé                                                  &  & \\
    \rh 60 & Mire de test (couleurs RGB, dégradé, damier)       & Affichage correct en 176$\times$176, couleurs justes, pas de pixel/colonne mort &  & \\
    \rh 61 & Orientation et origine des coordonnées             & Conformes au rendu attendu du framebuffer                                       &  & \\
    \rh 62 & Taux de rafraîchissement du framebuffer custom     & Fluide, pas de \textit{tearing} notable (à documenter)                          &  & \\
\end{longtable}

% =====================================================================
\section{Carte microSD (SPI3 + FATFS)}
% =====================================================================
\begin{longtable}{\protocolcols}
    \caption{Montage du système de fichiers et accès aux données.}
    \label{tab:proto-sd}                                                                                             \\
    \toprule \phdr \midrule \endfirsthead
    \multicolumn{5}{l}{\footnotesize\itshape Tableau \thetable{} — suite}                                            \\
    \toprule \phdr \midrule \endhead
    \midrule \multicolumn{5}{r}{\footnotesize\itshape suite page suivante}                                           \\ \endfoot
    \bottomrule \endlastfoot
    \rh 63 & Détection et montage FATFS                        & Montage réussi, capacité lue cohérente         &  & \\
    \rh 64 & Test écriture puis relecture d'un fichier         & Données relues identiques (comparaison/CRC OK) &  & \\
    \rh 65 & Lecture d'un fichier MP3 réel                     & Lecture des octets sans erreur                 &  & \\
    \rh 66 & Détection insertion/retrait (si broche CD câblée) & État cohérent                                  &  & \\
\end{longtable}

% =====================================================================
\section{Chaîne audio — signaux numériques (I\textsuperscript{2}S)}
% =====================================================================
\begin{longtable}{\protocolcols}
    \caption{Signaux d'horloge et de données I\textsuperscript{2}S (contenu à 44,1~kHz, 16~bits).}
    \label{tab:proto-i2s}                                                                                                                                                       \\
    \toprule \phdr \midrule \endfirsthead
    \multicolumn{5}{l}{\footnotesize\itshape Tableau \thetable{} — suite}                                                                                                       \\
    \toprule \phdr \midrule \endhead
    \midrule \multicolumn{5}{r}{\footnotesize\itshape suite page suivante}                                                                                                      \\ \endfoot
    \bottomrule \endlastfoot
    \rh 67 & Horloge BCK (oscilloscope) & Présente, $f =$ n\textsubscript{bits} $\times$ F\textsubscript{s} (p.~ex. 64 $\times$ 44,1~kHz $\approx$ 2,82~MHz selon config.) &  & \\
    \rh 68 & Horloge LRCK / WS          & $=$ F\textsubscript{s} (44,1~kHz)                                                                                                &  & \\
    \rh 69 & Données SDIN / DOUT        & Activité synchronisée avec LRCK                                                                                                  &  & \\
    \rh 70 & MCLK (si générée)          & Présente (p.~ex. 256 $\times$ F\textsubscript{s} $\approx$ 11,29~MHz selon config.)                                              &  & \\
\end{longtable}

% =====================================================================
\section{Chaîne audio — analogique}
% =====================================================================
\begin{longtable}{\protocolcols}
    \caption{Sortie DAC, filtre RC, amplificateur et jack (signal de test 1~kHz, volume numérique documenté).}
    \label{tab:proto-analog}                                                                                                                                                          \\
    \toprule \phdr \midrule \endfirsthead
    \multicolumn{5}{l}{\footnotesize\itshape Tableau \thetable{} — suite}                                                                                                             \\
    \toprule \phdr \midrule \endhead
    \midrule \multicolumn{5}{r}{\footnotesize\itshape suite page suivante}                                                                                                            \\ \endfoot
    \bottomrule \endlastfoot
    \rh 71 & Offset DC en sortie de jack (sans signal, ampli actif)  & $\approx$ 0~V (DirectDrive), qq mV max — \textbf{à vérifier avant de brancher un casque}                  &  & \\
    \rh 72 & Sortie différentielle PCM5242 (OUTL+/$-$, 1~kHz 0~dBFS) & Pleine échelle $\approx$ 4,2~V\textsubscript{rms} ; atténuer via le volume num. pour le test              &  & \\
    \rh 73 & Sortie après filtre RC (entrée ampli)                   & $\approx$ $-$6~dB vs sortie DAC (gain réseau 0,5), soit $\approx$ 2,1~V\textsubscript{rms} pleine échelle &  & \\
    \rh 74 & Sortie au jack (MAX97220, charge haute impédance)       & Plafond $\approx$ 2,15~V\textsubscript{rms} ; signal 1~kHz propre, non écrêté au volume nominal           &  & \\
    \rh 75 & Bruit de fond / souffle (sans signal, gain nominal)     & Faible, pas de sifflement ni bourdonnement audible                                                        &  & \\
    \rh 76 & Séparation des canaux : tonalité L seule                & Signal sur le canal gauche, droit silencieux (valide le mapping INL/INR)                                  &  & \\
    \rh 77 & Séparation des canaux : tonalité R seule                & Inverse du précédent                                                                                      &  & \\
    \rh 78 & THD+N à 1~kHz, niveau nominal (analyseur audio)         & Faible ($<$ 0,01~\% visé, à documenter)                                                                   &  & \\
    \rh 79 & SNR / plancher de bruit (analyseur audio)               & Élevé ($>$ 90--100~dB, à documenter)                                                                      &  & \\
    \rh 80 & Volume numérique PCM5242 (registres 61/62)              & Variation conforme (0,5~dB/pas), mute en butée                                                            &  & \\
    \rh 81 & Mapping potentiomètre $\rightarrow$ volume              & Cohérent et monotone du min au max                                                                        &  & \\
\end{longtable}

% =====================================================================
\section{Audio Bluetooth (A2DP)}
% =====================================================================
\begin{longtable}{\protocolcols}
    \caption{Liaison A2DP et contrôle AVRCP.}
    \label{tab:proto-bt}                                                                                                   \\
    \toprule \phdr \midrule \endfirsthead
    \multicolumn{5}{l}{\footnotesize\itshape Tableau \thetable{} — suite}                                                  \\
    \toprule \phdr \midrule \endhead
    \midrule \multicolumn{5}{r}{\footnotesize\itshape suite page suivante}                                                 \\ \endfoot
    \bottomrule \endlastfoot
    \rh 82 & Appairage avec un appareil source                       & Appairage réussi (codec SBC)                   &  & \\
    \rh 83 & Lecture en streaming                                    & Son présent au jack, sans coupure              &  & \\
    \rh 84 & Volume AVRCP (absolute volume)                          & Le réglage distant modifie le niveau de sortie &  & \\
    \rh 85 & Bascule filaire $\leftrightarrow$ Bluetooth (AudioSink) & Commutation propre, sans blocage ni glitch     &  & \\
\end{longtable}

% =====================================================================
\section{Tests fonctionnels et intégration}
% =====================================================================
\begin{longtable}{\protocolcols}
    \caption{Essais bout en bout et tenue en fonctionnement.}
    \label{tab:proto-func}                                                                                                                                            \\
    \toprule \phdr \midrule \endfirsthead
    \multicolumn{5}{l}{\footnotesize\itshape Tableau \thetable{} — suite}                                                                                             \\
    \toprule \phdr \midrule \endhead
    \midrule \multicolumn{5}{r}{\footnotesize\itshape suite page suivante}                                                                                            \\ \endfoot
    \bottomrule \endlastfoot
    \rh 86 & Lecture MP3 complète depuis microSD (bout en bout)                           & Décodage libhelix, audio continu, buffer jamais en sous-alimentation &  & \\
    \rh 87 & Navigation UI complète (boutons $\rightarrow$ écrans, push/pop)              & Réactive, pas de blocage de tâche                                    &  & \\
    \rh 88 & Autonomie sur batterie (durée à un volume donné)                             & À documenter ; cohérente avec la capacité de la cellule              &  & \\
    \rh 89 & Charge pendant la lecture (USB + lecture simultanés)                         & Fonctionne, la charge progresse, pas d'instabilité                   &  & \\
    \rh 90 & Comportement à la décharge profonde (VBAT bas)                               & Le XC6120N protège avant une tension dangereuse                      &  & \\
    \rh 91 & Thermique sous charge prolongée (TPS62A01, MAX77757, MAX97220, module ESP32) & Échauffement raisonnable, aucun point chaud anormal                  &  & \\
    \rh 92 & Stabilité longue durée (lecture prolongée)                                   & Stable : pas de reset, de brownout ni de fuite mémoire               &  & \\
\end{longtable}